\section{边警抵达内地以后：粮、马与州郡的选择}
\label{sec:eastern-logistics-interior}

边地告急抵达雒阳，并不会立刻变成一支已在路上的军队。粮要出仓，马要编配，
文书要经驿传转送，州郡还要在本地收成、役夫和守城之间决定先后。河西、凉州
和并州的军费因而不只是边疆问题：每一次长距离转运都会把内地的劳力、市场和
地方官的信誉接进防务。正史常保存将军的胜负，较少保存车马损耗、欠役和迁徙
家庭；这种沉默正是理解后期州郡为何更常自筹资源的起点。

西域都护与班超的使行也须从同一条物流线读。都护职名能协调使节和援军，不能
使绿洲、水源、向导和商队失去自己的计算。一次道路不通、一次驻点缺粮，便可
使朝廷的远方计划停在河西之外。甘英未至大秦不是可有可无的逸闻，而是提醒
读者：信息和外交要经过中间政权与海路，不能由皇帝的诏令直接校正。

黄巾和凉州动乱迫使地方同时面对军粮与治安。州牧、刺史和边将取得更多调度空间，
并非因为中央突然不再存在，而是共同的转运规则已无法及时覆盖多处危机。临时
能筹集粮饷的人获得了保护居民、任用部曲和同邻近势力谈判的机会；这种机会后来
才会成为区域军政的资源。边警向内地传递的，不只是敌情，也是一套重新分配谁能
掌握道路、仓储和人的政治后果。

\subsection{河西之外没有直线：使节的脚程与绿洲的选择}

一纸任命把\eventref{EH-0091-BAN-CHAO-PROTECTOR}{班超置为西域都护}，并不能把
雒阳同所有绿洲直接接在一起。都护的作用是协调使节、驻军与朝廷的关系；真正的
道路仍由泉水、城郭、商队、向导和各王廷的判断组成。使者离开河西以后，所带的
文书要在一站站的交接中保持可辨认，随行者要找到可以补给的地方，接待者也要
衡量同汉廷往来会给自己带来保护、贸易机会还是邻近势力的报复。把某国“服属”
写成地图上的一种颜色，会把这些反复发生的选择从历史中抹去。

这种距离也改变内地对边事的理解。一封从西域来的奏报抵达雒阳时，可能已经经过
数月的传递；中枢据此发布的指令再次抵达原处，还要面对季节、道路和人事是否
已经改变。信息延迟并不表示朝廷什么也不知道，而是说明它必须借助能够在路上
作判断的人。都护、属吏、商旅和当地首领并非诏令的透明管道，他们会解释消息、
寻找马匹、选择等待或前行。帝国的远方因此不是被命令征服的空白，而是一系列
需要持续维护的交接点。

\eventref{EH-0097-GAN-YING-ANXI}{甘英西使}尤其容易被后来关于“大秦”的好奇心
遮住。更值得停留的是，他到达安息属地而未能继续前行这一事实所暴露的条件。
海路信息、沿途接待与中间政权的利益不是汉廷能够单方规定的事。我们不宜替他
绘制一条精确而无争议的路线，也不能把后出的地理叙事当成行程日志；但可以看见，
一项远程外交必须把道路的风险交给许多不属于汉廷官僚体系的人。所谓“未至”
不是简单失败，它说明远方知识始终带着中介的边界。

这条边界与内地仓储有直接关系。西域来使不只是携带消息的人，也是要被供给、
安置和辨验的人；从内地向西送出的布帛、谷物、马匹和赏赐，同样需要在河西的
关塞和传舍接续。每增加一段距离，沿线的县乡便多一层可能被征用的车马与劳力。
朝廷不可能把所有费用逐户说明，地方也不可能只把来往看作遥远外交。车夫误期、
马匹病死、河谷封冻或一处仓粮不足，都可能把西域的政治关系转化为内地的差役
和价格问题。

\subsection{军粮不是一笔已经付清的数目}

边地军费最容易在史书中以一次“发兵”出现，仿佛粮食随命令同时出发。实际上，
军粮先要从某处收集，再要决定以谷、布、钱或牲畜中的哪一种计算，之后才有车、
舟、人和守卫把它送向下一站。沿路的县乡不能无限提供役夫：春耕、地方治安、
私人运输和别处的灾害都在争夺同一批劳力与畜力。一次转运若在中途变慢，远处
军营得到的不是抽象的“预算不足”，而是马匹少草、士卒少食、将领不得不缩减
行动或向附近征取。

这正是\eventref{EH-0107-LIANG-QIANG-WAR}{永初羌乱}为何既是边战又是内地问题。
陇西、金城及其相邻地区的冲突使行军、耕作和人户登记同时变得困难；迁出战区的
人又将安置和粮食压力带到新的郡县。我们无法从后出的正史叙事精确算出每一石
粮、每一名役夫或每一支部众的损失，因此不应把“耗费巨大”写成未经核对的数字。
但长期投入军费、反复迁徙和地方社会受压的关系是清楚的：当常规转运不能可靠
维持，能够就地筹粮、招募和守路的人便不再只是执行者。

这种地方筹给并非全然出于野心。边将若等待雒阳的每一次拨付，可能先失去城塞、
盟友或逃难者；州郡若把所有粮食送走，又可能无力应付本地治安。临时向富户借粮、
要求大族出人、改用附近市场或让军队自行取给，都是压力下可能出现的办法。它们
能够让一段道路暂时不断，却会把公共供给和私人保护缠得更紧。借出粮食者可能
获得声望、债权或任命，承担役使者则可能离开原有耕作；后来人们看到的“地方
势力”，往往正是在这种不易留名的交接中积累了可用资源。

汉廷并没有因此立即放弃对地方的要求。诏令、考课和任命仍试图把地方动作纳入
共同名义，地方官也仍需向上报告。变化在于，报告和供给不再总能由同一条稳定的
链条相连。消息可以先到，粮食可能后到；一支军队可以按朝命出发，却在半途改由
地方网络维持。把这一差距称为“中央失控”太粗糙，更好的说法是：国家赖以把
命令变成行动的物流接口，在不同地区以不同速度变得昂贵和脆弱。

\subsection{从凉州的车辙到许县的屯田}

一八四年凉州起事时，西北长期积累的军人、道路和供给关系进入了新的危机。
\eventref{EH-0184-LIANGZHOU-REBELLION}{凉州军民与羌胡部众的动乱}不应被理解为
一群统一的“边人”突然反对朝廷。不同军人、地方居民和被史书以羌胡概称的部众，
各自面对征发、牧地、交易、保护和首领关系的压力；记录更容易留下首领姓名，
较难留下这些关系怎样在具体城镇和山谷中变化。可确认的是，朝廷不得不把更大
的军政任务交给能够在当地取得人、马和粮的将领，因而使边地积累的资源获得进入
全国政治的通道。

后来曹操在许县经营屯田，所面对的也是供给怎样不再完全依赖即时掠取的问题，
只是场景从陇右的山道移到了战乱后的北方平原。屯田把流民、荒地、军队和仓储
连在一起，给部分家庭提供了重新耕作的路径，也把劳动、迁徙和军纪置于新的约束
之下。它不是凉州军费的简单解决，更不是所有地区可以照搬的答案；它显示区域
政权若想维持军队，仍须处理道路和粮食留下的旧问题。许县的诏令之所以能比别处
更有重量，不只是献帝在那里，而是人们试图用田、仓与官吏让这份名义能够反复
兑现。

从甘英止步的远方到凉州车队的迟滞，再到许县的屯田，物流并不是帝国政治背后
安静的技术层。它规定了谁可以等待、谁必须先行动，规定一项命令到达时是否还
有可用的人马，也规定地方社会把国家看作可预期的结算者还是遥远的索取者。汉末
的区域化并非由于道路突然消失；许多道路仍在使用。改变的是它们不再都把粮食、
信息和人稳定地送向同一个中心。
